\documentclass{templates/modern}
\hypersetup{
    pdftitle={Dupont's CV},
    pdfauthor={Dupont},
    pdfcreator={Dupont}
}

\definecolor{primaryColor}{RGB}{102, 126, 234}

\author{Jean Dupont}

\begin{document}
    \makeheader{}{Jean Dupont}{Paris, France}{jean.dupont@email.com}{+33 1 23 45 67 89}{https://jeandupont.dev}{https://linkedin.com/in/jeandupont}{https://github.com/jeandupont}

     \section{Profil}
        \begin{onecolentry}
        Développeur Full-Stack passionné avec 5 ans d'expérience dans la création d'applications web modernes. Expertise en JavaScript, React, Node.js et bases de données. Recherche de nouveaux défis techniques dans une équipe dynamique.        \end{onecolentry}
    \section{Formation}
    \experienceentry{2020}{}{Master en Informatique en Génie Logiciel}{Université Pierre et Marie Curie}{\item Mention Bien - Spécialisation en développement web et architecture logicielle
            }

    \section{Expérience}
    \experienceentry{Jan 2021 – Présent}{Paris, France}{Développeur Full-Stack Senior}{TechCorp Solutions}{\item Développement d'applications web modernes avec React et Node.js
            \item Gestion et encadrement d'une équipe de 3 développeurs juniors
            \item Optimisation des performances applicatives et renforcement de la sécurité
            \item Mise en place de pipelines CI/CD avec Docker et Kubernetes
            \item Conception d'architectures scalables et maintenables
            \item Participation aux décisions techniques et choix technologiques
            }

    \section{Projets}
    \projectentry{}{E-Commerce Platform}{\item Plateforme e-commerce complète développée avec React, Node.js, MongoDB et Stripe. Fonctionnalités: gestion des produits, panier, paiements sécurisés, dashboard admin.
            \item GitHub: \href{https://github.com/jeandupont/ecommerce-platform}{E-Commerce Platform}
            }

    \section{Certificats}
    \experienceentry{Mars 2023}{}{AWS Certified Developer}{Amazon Web Services}{\item Certification professionnelle en développement d'applications sur AWS
            \item Location: En ligne
            }

    \section{Compétences}
    \skillsentry{Langages de programmation}{JavaScript, TypeScript, Python, Java, PHP, SQL}

    \section{Langues}
\    \skillsentry{Langues}{Français (Native)}

\end{document}